% Section 9: approximate laws and the best system (owner: laws group).
% Plan: ../OUTLINE.md (sec:laws); rows: ../CLAIMS.md section 8; notation: ../NOTATION.md section 8.1.
% Full statements of the results stated here in words, and all proofs: app-laws.tex.
% Local symbols (renamed from the source items, NOTATION section 11 and local choices):
%   time-naming cost c_time (pi_c); measurement error zeta_j (eta_j); patch set Z (S); interval maps
%   \underline f, \overline f (l, u); completion distribution \mathcal G and band-position law \mathcal F
%   (G, F: theory and function machines are global); band position w (x); Euler's number \mathrm e;
%   per-observation literal cost gamma (c); module and theory lengths ||.|| without stop bit.
\section{Approximate laws and the best system}\label{sec:laws}

Kepler's laws are false. Against a modern planetary theory, a fixed ellipse misses Mars by one or two arcminutes over a few decades and the Moon by about a degree within a month (\cref{tab:kepler:residuals}); yet the laws were, and are, good science. A proof-requiring inducer seems ill suited to such laws: a deterministic theory either proves a datum or does not, and an exact law that is slightly wrong is refuted by any datum precise enough to show its error. This section asks what \SAIabbr does with an approximately true law, and what becomes, inside it, of Lewis's best system and of explanatory power.\footnote{The questions of approximately true laws and of the domain of applicability of a model are discussed in H\"anni's working notes \citep{hanni2026notes}.} In brief, the difficulty is a limitation of proof-requiring compatibility with deterministic theories, not of the axiom view as such: $\MA$ mixes randomised theories, and under \SAIabbr's pricing an approximately true law is a law with a calibrated error model (\cref{prop:laws:errormodel}). These results carry the row on approximate truth of the argument map (\cref{sec:phil:map}). \Cref{app:laws} states in full the results that this section states in words, and proves all of them.

\subsection{Setting}\label{sec:laws:setting}

The data are answers to questions (\cref{def:setting:questions}), such as the cell of width $\Delta$ in which the longitude of Mars lies at time $t$ (\cref{def:kepler:language}).

\begin{definition}[Answered data, exception codes, approximate laws]\status{definition}\src{laws Defs 1.1--1.3, 1.5, 1.7; \S2 setting; Setting S3}\label{def:laws:setting}
\leavevmode
\begin{enumerate}[label=(\alph*),itemsep=1pt,topsep=2pt]
\item \emph{Answered data.} A datum is a pair $(\mathcal Q,a)$ of a question and its true answer; a data set $D$ has pairwise distinct questions and is read as the two-valued data set $\{(a,1):(\mathcal Q,a)\in D\}$; $\dom D$ lists its questions, $y_D$ their answers, and $n:=\lvert D\rvert$. A theory $T$ \emph{proves} (\emph{refutes}) the datum if $B\cup T\vdash a$ ($B\cup T\vdash\neg a$), and is \emph{silent} otherwise; $A_T(\mathcal Q):=\{a'\in\mathcal Q:B\cup T\nvdash\neg a'\}$, and $\gamma(\mathcal Q):=\log\lvert\mathcal Q\rvert$ is the \emph{literal cost}.
\item \emph{Codes.} For a prefix-free code $c_S$ of sentences, $\ell(\chi):=1+\lvert c_S(\chi)\rvert$, and $\lVert\Gamma\rVert:=\sum_{\chi\in\Gamma}\ell(\chi)$ for a finite set $\Gamma$ of sentences (a theory or a \emph{module} of one). If $B\cup T\cup D^\pm$ is consistent, the \emph{exceptions} of $T$ are the data it does not prove, priced by $\Exc^Q_D(T):=\sum\gamma(\mathcal Q)$ (code Q), $\Exc^F_D(T):=\sum\ell(a)$ (code F) or $\Exc^R_D(T):=-\log P^R_T(D)$ (code R, below); otherwise $\Exc_D(T):=\infty$. Unmarked, $\Exc$ is code Q. $J_\lambda(T;D):=\lambda\lVert T\rVert+\Exc_D(T)$ for $\lambda>0$ (and $J^Q_\lambda$, $J^F_\lambda$, $J^R_\lambda$); its minimisers are the \emph{$\lambda$-MAPs} \citep[Definition~3.6]{claudeSFI2026steeper}. $T^{-Z}$ withdraws $T$ at the finitely many times in $Z$, at about $c_{\rm time}$ bits per time.
\item \emph{Interval laws.} $B:=\mathrm{RCF}$ (decidable and complete; standard, \citealp{tarski1951decision}), with constants $q_t$ for times $t$ in a decidable set $I$ of numerals, and a rational range $[0,\mathsf R]$. A \emph{pointwise interval law} is $\{\underline f(t)\le q_t\le\overline f(t):t\in I\}$ for total computable $\underline f\le\overline f:I\to\mathbb Q\cap[0,\mathsf R]$ named by a codeword of a fixed decidable prefix-free presentation language; $\lVert T\rVert$ is the codeword's length. The \emph{exact law} $T_0$ has $\underline f=\overline f=f$, the \emph{honest law} $T^\varepsilon$ has $f\mp\varepsilon$ (clipped). \emph{Interval data} answer questions that partition $[0,\mathsf R)$ into finitely many cells with rational endpoints; with $C_t$ the intersection of the cells that $D$ assigns to $t$, $P^R_T(D):=\prod_t\lvert C_t\cap[\underline f(t),\overline f(t)]\rvert/(\overline f(t)-\underline f(t))$ (length; the factor is $\ind[\underline f(t)\in C_t]$ if $\underline f(t)=\overline f(t)$).
\item \emph{Errors.} For a real quantity $q(t)$, $e(t):=f(t)-q(t)$, and $\varepsilon^*:=\operatorname{ess\,sup}\lvert e\rvert$ under the law of the error at a data time. \emph{Regime G:} datum $j$ is ``$q(t_j)\in[v_j-\delta,v_j+\delta]$'' with $v_j=q(t_j)+\zeta_j$ rational and $\lvert\zeta_j\rvert\le\delta_0\le\delta$. \emph{Regime R:} $q$ lives on a circle of circumference $\mathsf R$; datum $j$ is about a \emph{fresh} time $t_j$ (an injective computable sequence of numerals); the errors $e_j:=e(t_j)$ are i.i.d.; and datum $j$ is the cell of width $\Delta$ ($\mathsf R/\Delta$ an integer) containing $q(t_j)$ in a grid shifted by a dither $u_j$, uniform and independent of everything else; $\gamma=\log(\mathsf R/\Delta)$, $p_\Delta:=\Ex\min(\lvert e\rvert/\Delta,1)$ and $s_\Delta:=\min(2\varepsilon^*/\Delta,1)$.
\item \emph{Peakedness.} $\pk(e):=\varepsilon^*/\Ex\lvert e\rvert$ and $\pk_\varepsilon(e):=\varepsilon/\Ex\lvert e\rvert$; every reported value names its fit (minimax or mean).
\end{enumerate}
\end{definition}

\subsection{Deterministic theories and precision}\label{sec:laws:strength}

Over deterministic theories, proof-requiring compatibility is monotone in logical strength, so data never move the odds towards a weaker theory: $T_0\vdash T^\varepsilon$, every data set compatible with $T^\varepsilon$ is compatible with $T_0$, and with prior weights $2^{-\lVert T\rVert}$ the posterior odds of $T_0$ against $T^\varepsilon$ are at least $2$ on each, although $T_0$ is false when $e\not\equiv0$ (\cref{prop:laws:strength}). Under guaranteed tolerances (regime G), $T_0$ is almost surely compatible with the data at every sample size when $\delta\ge\varepsilon^*+\delta_0$ and, under mild conditions on the measurement error, is eventually refuted when $\delta<\varepsilon^*+\delta_0$ (\cref{prop:laws:tolerance}). Rounding changes the picture.

\begin{proposition}[Dithered rounding]\status{proved; computed}\src{laws Prop 2.5}\label{prop:laws:rounding}
In regime R, given $e(t_j)$, $T_0$ proves datum $j$ with probability $(1-\lvert e(t_j)\rvert/\Delta)^+$ and refutes it otherwise. So $T_0$ survives $n$ data with probability $(1-p_\Delta)^n$, its first refutation has a geometric index with mean $1/p_\Delta$, and it is almost surely refuted eventually iff $\Prob[e\ne0]>0$, for every $\Delta$, however coarse. For $\Delta\ge\varepsilon^*$, $p_\Delta=\Ex\lvert e\rvert/\Delta$: the exact law survives about $\Delta/\Ex\lvert e\rvert$ data.
\end{proposition}

This corrects the natural expectation that coarse data protect a false exact law: precision belongs to the data set, not to the datum, and an exact law with a sinusoidal error of amplitude $1'$ is refuted by positions rounded to a dithered grid of $1^\circ$ after about $94$ data (the proof and the comparison of several exact laws, \cref{prop:laws:choice}, are in \cref{app:laws:precision}). We have not found these two regimes stated elsewhere; precision-dependent selection as such is familiar from accounts based on AIC, MDL and MML (cf.\ \citealp{forster1994how,rissanen1978modeling,wallace2005statistical}).

\subsection{Randomised theories and laws with error models}\label{sec:laws:random}

A deterministic theory pays, in code F, for which questions are exceptions (\cref{prop:laws:exceptions}). A theory machine that reads a random tape answers every question at once, at random inside what the law allows: for a pointwise interval law $T$, $\LA(D)\le\lVert T\rVert+\Exc^R_D(T)+c$ on interval data consistent with $T$ (\cref{prop:laws:randomised}). In fact $\LA(D)=\min_G[K(G)-\log\nu_G(D)]\pm c$ over theory machines $G$: \SAIabbr is the best two-part code over randomised theories, neither $\KA$ nor literal-exception codes over finite theories are within an additive constant of it in general, and over all theory machines a steeper penalty collapses onto the universal machine, so best-system analyses use laws with private literal facts (\cref{prop:laws:twopart,prop:laws:P2hyp}). \emph{\SAIabbr's pricing} of a form is the weight $2^{-(K(G)-\log\nu_G(D))}$ of a named theory machine $G$ in $\MA$'s mixture, not the whole posterior.

With exceptions named in the law at $c_{\rm time}$ bits each (a cost model valid only while times are that cheap to name; \cref{prop:laws:patches}), the honest law $T^{\varepsilon^*}$ has a lower rate than the patched exact law $T_0^{-Z}$, for $\Delta\ge2\varepsilon^*$, iff $\lambda c_{\rm time}>\gamma(2\pk(e)-1)$: spread errors favour honesty and peaked ones exceptions (\cref{prop:laws:peaked,prop:laws:quantile}). Under random completion silence costs at most $\log\mathrm e$ bits per straddled cell boundary, $\mathrm e$ Euler's number (\cref{lem:laws:codeR}), and the comparison moves inside the band. For $\varepsilon\ge\varepsilon^*$ and a computable distribution $\mathcal G$ on $[0,1]$, the \emph{band-confined completion} $\mathrm{BC}_{\mathcal G}$ emits $T^\varepsilon$ and, for each time $t_i$, the value $f(t_i)-\varepsilon+2\varepsilon\hat w_i$ with $\hat w_i\sim\mathcal G$ drawn from the tape. On data at fresh times its law is $\prod_j\mathcal G(C'_j)$, with $C'_j$ the cell of datum $j$ in band coordinates (the band mapped affinely onto $[0,1]$), and $K(\mathrm{BC}_{\mathcal G})\le\lVert T^\varepsilon\rVert+K(\mathcal G)+c$. Uniform $\mathcal G$ gives the \emph{uniform band}; $\mathrm{HY}(\rho)$, with $\mathcal G=(1-\rho)\delta_{1/2}+\rho\,\mathrm{Unif}[0,1]$, is ``exact, except in a share $\rho$ of cases, and then within $\varepsilon$''.

\begin{proposition}[Laws with error models]\status{proved; computed}\src{laws Prop 3.6a, Def 3.6a}\label{prop:laws:errormodel}
In regime R let $\varepsilon\ge\varepsilon^*$, and let $\mathcal F$ be the distribution function of $w_j:=(\varepsilon-e_j)/(2\varepsilon)$, the position of the true value in the band.
\begin{enumerate}[label=(\alph*),itemsep=1pt,topsep=2pt]
\item If $\Delta\ge2\varepsilon$, the rate of $\mathrm{BC}_{\mathcal G}$ is $(2\varepsilon/\Delta)\,c_{\mathcal G}$, with
\[c_{\mathcal G}:=\Ex_w\int_0^1\big(\ind[v<w](-\log(1-\mathcal G(v)))+\ind[v>w](-\log\mathcal G(v))\big)\,dv.\]
So band-confined completions are ranked in the same way at every such $\Delta$.
\item $c_{\mathcal G}\ge c_{\mathcal F}:=\int_0^1H_2(\mathcal F(v))\,dv$, $H_2$ the binary entropy, with equality iff $\mathcal G=\mathcal F$ almost everywhere on $[0,1]$. In particular the uniform band is the best band-confined completion iff $e$ is uniform on $[-\varepsilon,\varepsilon]$.
\item For every theory machine $G'$, $G_U$ included, almost surely $\liminf_n\frac1n(-\log\nu_{G'}(D_n))\ge r_{\mathcal F}(\Delta):=\Ex[-\log\mathcal F(C'_1)]$, the conditional entropy rate of the answers given the questions, while $-\frac1n\log\nu_{\mathrm{BC}_{\mathcal F}}(D_n)\to r_{\mathcal F}(\Delta)$. If $\mathcal F$ is computable, then $\LA(D_n)/n\to r_{\mathcal F}(\Delta)$ almost surely. For $\Delta\ge2\varepsilon$, $r_{\mathcal F}(\Delta)=(2\varepsilon/\Delta)\,c_{\mathcal F}$.
\item $c_{\mathrm{HY}(\rho)}$ is convex in $\rho\in(0,1]$, and some $\mathrm{HY}(\rho)$ with $\rho<1$ beats the uniform band iff $\Ex\ln(1+\lvert e\rvert/\varepsilon)<2\ln2-1$. This requires $\pk_\varepsilon(e)>\ln2/(2\ln2-1)\approx1.79$, and holds whenever $\pk_\varepsilon(e)>\mathrm e/(4-\mathrm e)\approx2.12$.
\end{enumerate}
\end{proposition}

\emph{Proof idea.} A cell boundary falls into the band with probability $2\varepsilon/\Delta$, at a uniform position $v$, and the datum then costs $-\log$ of the $\mathcal G$-mass on the true side; (b) is Gibbs' inequality given $v$; (c) uses that a theory proves at most one answer list, with Markov's inequality and the Borel--Cantelli lemma; (d) is convexity and the sign of the derivative at the uniform band (\cref{app:laws:honest}).

So under \SAIabbr's pricing an approximately true law is a law with a calibrated error model, not a silent partial theory and not an exact law with listed exceptions; the dichotomy between the two survives inside the band, with peakedness thresholds near $2$. That a model should include its noise distribution is the standard situation of MDL (cf.\ \citealp{rissanen1978modeling,grunwald2007minimum}); what is specific is the reading for a proof-requiring inducer, in which the honest partial theory is one completion among many. Selection among finitely many forms is almost sure (\cref{thm:laws:as,cor:laws:phases}). For deterministic theories the trichotomy is the formal home of humility: on a finer question, the honest and the empty law put their posterior mass on ``independent'', and whether the inducer ends humble is decided not by the prior, which favours the exact law, but by the error's shape and the price of silence (\cref{prop:laws:humility}). For the randomised theories that $\MA$ mixes, humility is calibration: by (c) the preferred spread over the cells is the true error distribution.

\paragraph{Kepler's residuals.} For seven bodies and six windows each (three for the Moon) we kept the fixed Kepler ellipse with the least largest longitude error (the minimax fit) and the one with the least mean error (the mean fit), and took their errors at $8000$ uniform times as error laws, against the plan94 and moon98 ephemerides (\cref{tab:laws:residuals} in \cref{app:laws:honest}). The errors are moderately spread: the peakedness is $1.60$--$3.95$ for the minimax fit and $1.89$--$7.08$ for the mean fit, against $\pi/2$ for a sinusoid and $2$ for uniform $\lvert e\rvert$. A computation of the rates at eleven cell widths from $0.1'$ to $30^\circ$ (code in the accompanying repository; details in \cref{app:laws:honest}) gives the following. Among the Kepler-based theory machines compared, the law with a 40-bin error histogram, fitted on half of the times and scored on the other half, beats the uniform band and the exact law with exceptions in all $429$ (body, window, cell width) cells, by $0.4\%$ to $25.7\%$ of the band's rate, although the uniform band wins every cell within the pool $\{\emptyset,T_0^{-Z},\text{uniform band}\}$; and in the coarse regime the two-point form beats the band for $34$ of the $39$ body--window pairs, exactly where \cref{prop:laws:errormodel}(d) says. The residuals are deterministic functions of time, so a better theory (Newton's, or the ephemeris itself) eventually beats every Kepler-based machine.

\subsection{Sound reasoning with approximate models}\label{sec:laws:sound}

If the background is true, true honest models of one quantity are jointly consistent, and every statement about the quantity derived from them is a set containing its true value. If the background also decides the comparisons involved (as $\mathrm{RCF}$ does for semialgebraic models), exact approximations that differ anywhere on a common domain are jointly inconsistent and so prove everything, while joint consistency of their honest forms is a pairwise check of tolerances in one dimension and for boxes, and an $(m+1)$-wise check for Euclidean balls in $\mathbb R^m$, by Helly's theorem (\cref{prop:laws:sound}). An honest dynamical law with initial tolerance $\varepsilon^{(0)}$ and tolerance $\varepsilon$ per step proves, after $r$ steps of an $L$-Lipschitz map, the tolerance $L^r\varepsilon^{(0)}+\varepsilon(L^r-1)/(L-1)$ ($L\ne1$), tightly for linear maps in one dimension: it has a derivation horizon, and an exact approximation is guaranteed correct up to that horizon (for linear maps in one dimension, only up to it) (\cref{prop:laws:horizon,prop:laws:exacthorizon}). The device is sound, but it does not say how to choose a model, a domain or a tolerance (\cref{sec:limits:domains}).

\begin{remark}[Typical against worst case]\status{computed}\src{laws Rem 4.4}\label{rem:laws:typical}
With $L=1$ and $r=100$ steps the propagated tolerance is $100\varepsilon$, while with independent random-sign errors $\pm\varepsilon$ the mean deviation is $\varepsilon r\binom{r}{r/2}2^{-r}\approx7.96\varepsilon$: a probabilistic error model answers much further out, with a stated confidence instead of a guarantee.
\end{remark}

\subsection{The best system as a steeper MAP}\label{sec:laws:best}

On Lewis's account the laws are the theorems of the deductive system that best balances simplicity against strength (cf.\ \citealp{lewis1983new,lewis1994humean}). Take strength to be fit to data in a fixed base language, the balance to be $\lambda$, and the candidates a class of laws with private literal facts; the relation to Lewis is an analogy, since best systems are a metaphysics of laws and \SAIabbr is inference from finite data. The points $(\lVert T\rVert,\Exc_D(T))$ have a convex hull with vertices $v_i$ in order of length, joined by edges of slopes $-\sigma_i$, $\infty=:\sigma_0>\sigma_1>\dots>0$; for $\lambda\in(\sigma_{i+1},\sigma_i)$ the $\lambda$-MAPs are the theories at $v_i$ \citep[Lemma~3.8]{claudeSFI2026steeper}. A point $(\alpha_{\rm r},L_{\rm r})$ attained by some theory is a \emph{zero-window $(\sigma,\tau)$-kink} \citep[Definition~5.2]{claudeSFI2026steeper} if every theory shorter by $x>0$ has fit at least $L_{\rm r}+\sigma x$ and every theory longer by $x>0$ has fit at least $L_{\rm r}-\tau x$.

\begin{proposition}[Kind nature is a kink]\status{proved, given Lemma 3.8 and Corollary 5.4 of the companion paper on steeper priors; computed}\src{laws Prop 5.6}\label{prop:laws:kink}
For a family of candidate theories satisfying the hypotheses of the hull lemma (\cref{prop:laws:P2hyp}) and $0\le\tau<\sigma\le\infty$, the following are equivalent: (a) the set of $\lambda$-MAPs is the same for every $\lambda\in(\tau,\sigma)$; (b) some hull vertex $v_i$ has $\sigma_{i+1}\le\tau$ and $\sigma\le\sigma_i$; (c) the point family has a zero-window $(\sigma,\tau)$-kink at a point that is a hull vertex.
\end{proposition}

\emph{Proof idea.} The $\lambda$-MAP set is fixed between consecutive edge slopes and changes at each; (b)$\Leftrightarrow$(c) is \citet[Corollary~5.4(c)]{claudeSFI2026steeper} (\cref{app:laws:best}).

So nature is kind on an interval of balances exactly when the data's hull has a kink there. In a solar-system toy with eight planets of $655$ to $3838$ observations, sixty asteroids with a median of about five and thirty comets with one or two (\cref{prop:laws:K3}), the $\lambda$-MAP is ``the law plus the elements of all eight planets'' for every $\lambda\in(7.03,135.37)$, while below $\lambda\approx7$ it changes with every asteroid (code in the accompanying repository). Inside such an interval the worry that the best system depends on an arbitrary choice of balance (cf.\ \citealp{woodward2014simplicity}) is moot, and kindness is a checkable property of a language, a family of candidates and a data set. Lewis's trivialisation worry, that a predicate true of exactly the actual things makes its universal closure maximally simple and strong (cf.\ \citealp{lewis1983new}), is blocked by the coding bound, up to the complexity of a language chosen after the data (\cref{prop:laws:trivial}).

\begin{definition}[Explanatory power]\status{definition}\src{laws Def 5.2}\label{def:laws:EP}
For a theory $T$ and a module $\mathcal A\subseteq T$, $\EP_D(\mathcal A;T):=\big(\Exc_D(T\setminus\mathcal A)-\Exc_D(T)\big)/\lVert\mathcal A\rVert$: the bits of data the module explains per bit of its statement, in the context $T$.
\end{definition}

Every module of a $\lambda$-MAP has $\EP\ge\lambda$, no addition to it has $\EP>\lambda$, and for $\lambda>1$ a fact used once never enters (\cref{prop:laws:slopes,prop:laws:once}). So the model--randomness split of the companion paper separates facts reused across many data from facts used once, not laws from initial conditions: in the toy, the elements of a well-observed body explain hundreds of bits per bit and sit on the steep side of the hull with the law (\cref{prop:laws:K3}). What singles out the law is how its $\EP$ scales: with $m_{\rm b}$ bodies of $n_0$ observations each, all stated, a law of $L$ bits has $\EP=m_{\rm b}n_0\gamma/L$, which grows with the number of systems, while each body's elements explain only that body's observations. Kepler's third law $\mathrm{K3}$ is a law about initial conditions: it shortens every elements module by $s$ bits, and when the set $S$ of bodies whose elements are stated is the same with and without it, it lowers the objective iff $s\lvert S\rvert$ exceeds its length, whatever $\lambda$ is; this compares law choices at exchange rate $1$ and is not membership of the $\lambda$-MAP. In the question format, elements are question antecedents and never enter a $\lambda$-MAP (\cref{prop:laws:questions}). \Cref{sec:kepler} computes what $\mathrm{K3}$ is worth in Kepler's setting.

\subsection{Explanation and correction}\label{sec:laws:explain}

$\lVert\mathcal A\rVert\,\EP_D(\mathcal A;T)$ is the log-likelihood ratio of $T$ against $T\setminus\mathcal A$ for $P_T(D):=2^{-\Exc_D(T)}$. The support a hypothesis gets from joint data $D_1\cup D_2$ splits into the supports from the parts plus a term $U$, the change in the mutual information between $D_1$ and $D_2$ that the hypothesis brings about, which has the form of Myrvold's measure of unification \citep{myrvold2003bayesian}: with stated parameters, for laws that are products over questions, $U=0$, and when a $k$-bit parameter is integrated out, $U$ is $k$ within one system, $0$ across systems with independent parameters, and the bits of a shared constant across systems that share one (\cref{prop:laws:EP}). At $\lambda=1$ a module pays iff its $\EP$ exceeds $1$, the compression criterion of inductive logic programming (cf.\ \citealp{muggleton1995inverse,srinivasan1994justification}). Copernicus against Ptolemy is unification without better prediction: against Ptolemy with the coincidence of each superior planet's retrograde motion and opposition stated as an axiom, at $a_p$ bits each, which predicts the same, Copernicus, at $H$ extra law bits, wins at every $\lambda$ iff $H<\sum_pa_p$, and against Ptolemy with free phases it wins iff $\lambda<\sum_pa_p/H$ (\cref{prop:laws:copernicus}; the lengths are not measured; for the history, see \citealp{kuhn1957copernican}).

Newton's theory does not entail Kepler's laws; it contradicts them (cf.\ \citealp{duhem1954aim}) and derives corrected statements, as in reduction to a corrected reduced theory (cf.\ \citealp{schaffner1967approaches}). In an exact rational toy, uniform circular motion perturbed to $x_{t+1}=Ax_t$ with $A$ a rotation plus $\eta_{\rm c}\mathbf P$, the successor proves the honest old law on the window $t\le\bar t$ with uniform tolerance $\varepsilon$ iff $\varepsilon\ge\varepsilon_{\min}(\bar t)$, where $\varepsilon_{\min}(\bar t)\le((1+\eta_{\rm c}\lVert\mathbf P\rVert)^{\bar t}-1)\lvert a\rvert$ for the initial state $a$, and, if $\eta_{\rm c}\mathbf Pa\ne0$, it refutes the exact old law; among uniform-tolerance forms, the strongest it proves is the strongest both theories prove, while time-dependent tolerances are stronger common consequences and the whole common content is the disjunction of the two theories (\cref{prop:laws:correction}). The old law's domain of applicability is derived, not stipulated. Deriving ``Kepler within $\varepsilon$ on a window'' from Newton needs validated perturbation bounds (\cref{sec:kepler:newton}).

\subsection{Domains of applicability and retention}\label{sec:laws:domains}

A domain restriction that leaves one system's data literal pays asymptotically iff the law, even with an honest tolerance for that system, is worse than literal data, exact ties being decided by the constants (\cref{prop:laws:domain}). Under \SAIabbr's pricing that never happens, at cell widths below $\mathsf R/\mathrm e$, for a system that the law confines within about $66^\circ$ of longitude (a sufficient bound): an honest band of half-width $\varepsilon^*$ costs at most $\log(2\varepsilon^*/\Delta)+\log\mathrm e$ bits per datum, below the literal $\log(\mathsf R/\Delta)$ iff $\varepsilon^*<\mathsf R/(2\mathrm e)$. So the Moon ($\varepsilon^*=2.6^\circ$ over a year) is kept in Kepler's law with its own tolerance at every cell width computed, from $10'$ to $30^\circ$, whereas under literal pricing at $10'$ and $60'$ a cheap domain clause excludes it. A superseded theory is retained where its derivation savings exceed its loss of fit, for all cell widths above a computable $\Delta^*$ within the range where its rate is monotone (\cref{prop:laws:retention}).

\paragraph{Domains are learned as tolerances.}\label{rem:laws:domainlearning}\status{assessment}\src{laws \S8}
We conclude that under \SAIabbr's pricing a domain of applicability shows up as a tolerance and an error distribution that depend on the system and the window, not as a sharp boundary: for every planet and window computed, and for the Moon over a month, a year and $18.6$ years, a law with its own tolerance has a lower rate than leaving the body's data literal. A sharp boundary is selected only when even an honest tolerance is worse than literal data, or under literal pricing at fine precision. The posterior over an explicit pool of domain predicates is not computed (\cref{q:laws:domains}).
